% Sections 16.1 to 16.5, translated from sv/chapters/16/theory.tex.

\section{Euler's formula}
\label{c16:sec:euler}
\textbf{Euler's formula} connects complex exponents and trigonometry:
\[
  e^{j\delta} = \cos\delta + j\sin\delta
\]
This means that every complex number can be written in \textbf{Euler form}:
\[
  z = \abs{z} \cdot e^{j\delta}
\]
Some important special cases are
\begin{gather*}
  e^{j\pi} = -1 \quad \text{(Euler's identity)}, \\
  e^{j\pi/2} = j, \qquad e^{j0} = 1, \qquad e^{-j\delta} = \cos\delta - j\sin\delta.
\end{gather*}

\section{Three forms of complex numbers}
\label{c16:sec:former}
A complex number can be written in three ways, and each form has its strengths.

\begin{booktable}{@{}p{7em}p{6em}L@{}}
  \thead{Form} & \thead{Notation} & \thead{Comment} \\
  \midrule
  Rectangular & $x + jy$ & Best for addition and subtraction \\
  Polar & $\abs{z}\,\angle\,\delta$ & Clear geometric interpretation \\
  Euler form & $\abs{z}e^{j\delta}$ & Best for multiplication and division \\
\end{booktable}

\noindent\textbf{Multiplication} in polar form:
\[
  z_1 \cdot z_2 = \abs{z_1}\abs{z_2}\,\angle\,(\delta_1 + \delta_2)
\]
\textbf{Division} in polar form:
\[
  \frac{z_1}{z_2} = \frac{\abs{z_1}}{\abs{z_2}}\,\angle\,(\delta_1 - \delta_2)
\]

\section{Phasors}
\label{c16:sec:fasor}
A \textbf{phasor} is a complex number that represents a sinusoidal signal. A voltage
$u(t) = \abs{U}\sin(\omega t + \delta)$ is represented by the phasor
\[
  U = \abs{U}\,\angle\,\delta = \abs{U}e^{j\delta}.
\]
Phasors simplify calculations with sine waves in electric circuits: adding sine waves reduces to
adding complex numbers.

\section{Converting between Euler form and the time domain}
\label{c16:sec:omskrivning}
From phasor to time domain:
\[
  U = \abs{U}e^{j\delta} \quad \Rightarrow \quad u(t) = \abs{U}\sin(\omega t + \delta)
\]
From time domain to phasor:
\[
  u(t) = \abs{U}\sin(\omega t + \delta) \quad \Rightarrow \quad U = \abs{U}\,\angle\,\delta
\]

\needspace{12\baselineskip}
\section{Worked examples}
\label{c16:sec:typexempel}

\begin{exempel}[Vectors as complex numbers]
The vectors $\mathbf{a} = (3, 4)$ and $\mathbf{b} = (-2, 4)$ are written as complex numbers and
added:
\begin{gather*}
  a = 3 + j4, \qquad b = -2 + j4, \\
  a + b = (3 - 2) + j(4 + 4) = 1 + j8, \\
  \abs{a + b} = \sqrt{1 + 64} = \sqrt{65} \approx 8.06.
\end{gather*}
\end{exempel}

\begin{exempel}[Euler form to rectangular form]
Write $u(t) = 3e^{j(100\pi t - \pi/4)}\,\text{V}$ in rectangular form at $t = 0$.
\begin{align*}
  u(0) &= 3e^{-j\pi/4}
        = 3\left(\cos\left(-\frac{\pi}{4}\right) + j\sin\left(-\frac{\pi}{4}\right)\right) \\
       &= 3\left(\frac{\sqrt{2}}{2} - j\frac{\sqrt{2}}{2}\right) \approx 2.12 - j2.12\,\text{V}
\end{align*}
\end{exempel}

\begin{exempel}[Phasor to time domain]
The phasor $I = 5 - j4\,\text{mA}$ belongs to the angular frequency $\omega = 100\pi\,\text{rad/s}$.
Write $i(t)$ in Euler form.
\begin{gather*}
  \abs{I} = \sqrt{25 + 16} = \sqrt{41} \approx 6.40\,\text{mA}, \\
  \delta = \arctan\left(\frac{-4}{5}\right) \approx -38.7^{\circ} \approx -0.675\,\text{rad}, \\
  i(t) = 6.40 \cdot e^{j(100\pi t - 0.675)}\,\text{mA}.
\end{gather*}
\end{exempel}
